\chapter{Évaluation du modèle}

\section{Introduction}
Dans ce chapitre, nous présentons l'évaluation des deux modèles proposés pour la détection
des textes générés par intelligence artificielle : un modèle LSTM (\textit{Long Short-Term Memory})
et un modèle CNN (\textit{Convolutional Neural Network}). L'objectif est de mesurer les performances
de chacun à l'aide de différentes métriques et d'analyser les résultats obtenus.

\section{Métriques d'évaluation}

Pour évaluer les performances des modèles, plusieurs métriques complémentaires sont utilisées.

\textbf{L'accuracy} (précision globale) mesure la proportion de prédictions correctes parmi
l'ensemble des prédictions réalisées :
\[
\text{Accuracy} = \frac{TP + TN}{TP + TN + FP + FN}
\]

\textbf{La précision} (\textit{Precision}) indique la proportion de textes correctement identifiés
comme générés par IA parmi tous ceux prédits comme tels, permettant de mesurer le taux de faux positifs :
\[
\text{Precision} = \frac{TP}{TP + FP}
\]

\textbf{Le rappel} (\textit{Recall}) mesure la capacité du modèle à détecter l'ensemble des textes
générés par IA présents dans le jeu de données :
\[
\text{Recall} = \frac{TP}{TP + FN}
\]

\textbf{Le F1-score} combine la précision et le rappel afin de fournir une mesure équilibrée des
performances, particulièrement utile en cas de déséquilibre entre les classes :
\[
\text{F1-score} = 2 \times \frac{\text{Precision} \times \text{Recall}}{\text{Precision} + \text{Recall}}
\]

où $TP$, $TN$, $FP$ et $FN$ désignent respectivement les vrais positifs, vrais négatifs,
faux positifs et faux négatifs.

\section{Résultats expérimentaux}

Les deux modèles ont été entraînés sur 75\,000 exemples, validés sur 10\,000 exemples et
évalués sur 15\,000 exemples. L'entraînement a été effectué sur 5 époques avec l'optimiseur
Adam et la fonction de perte \textit{Binary Cross-Entropy} (\texttt{BCELoss}).

\subsection{Résultats du modèle LSTM}

Le modèle LSTM est composé d'une couche d'embedding, d'une couche LSTM unidirectionnelle,
d'une couche de dropout ($p=0{,}5$) et d'une couche fully connected. Le tableau~\ref{tab:lstm_results}
présente les résultats obtenus sur le jeu de test.

\begin{table}[h]
\centering
\label{tab:lstm_results}
\begin{tabular}{lcccc}
\hline
\textbf{Classe} & \textbf{Précision} & \textbf{Rappel} & \textbf{F1-score} & \textbf{Support} \\
\hline
Humain (0)      & 0.95 & 0.97 & 0.96 & 7\,450 \\
IA (1)          & 0.97 & 0.95 & 0.96 & 7\,550 \\
\hline
\textbf{Accuracy globale} & \multicolumn{4}{c}{\textbf{0.96}} \\
\textbf{Macro avg}        & 0.96 & 0.96 & 0.96 & 15\,000 \\
\textbf{Weighted avg}     & 0.96 & 0.96 & 0.96 & 15\,000 \\
\hline
\end{tabular}
\caption{Rapport de classification -- Modèle LSTM}
\end{table}

La perte sur le jeu de test (\textit{Test Loss}) s'établit à \textbf{0.127}, et l'accuracy de test
atteint \textbf{96\%}. La matrice de confusion du modèle LSTM est représentée à la
figure~\ref{fig:lstm_confusion}.

\subsection{Résultats du modèle CNN}

Le modèle CNN est composé d'une couche d'embedding, de trois couches convolutionnelles 1D
avec des filtres de tailles $\{3, 4, 5\}$ et de nombres de filtres $\{100, 200, 300\}$,
d'une opération de max-pooling, d'une couche de dropout, et d'une couche fully connected.
Le tableau~\ref{tab:cnn_results} présente les résultats obtenus.

\begin{table}[h]
\centering
\label{tab:cnn_results}
\begin{tabular}{lcccc}
\hline
\textbf{Classe} & \textbf{Précision} & \textbf{Rappel} & \textbf{F1-score} & \textbf{Support} \\
\hline
Humain (0)      & 1.00 & 1.00 & 1.00 & 7\,450 \\
IA (1)          & 1.00 & 1.00 & 1.00 & 7\,550 \\
\hline
\textbf{Accuracy globale} & \multicolumn{4}{c}{\textbf{0.998}} \\
\textbf{Macro avg}        & 1.00 & 1.00 & 1.00 & 15\,000 \\
\textbf{Weighted avg}     & 1.00 & 1.00 & 1.00 & 15\,000 \\
\hline
\end{tabular}
\caption{Rapport de classification -- Modèle CNN}
\end{table}

La perte sur le jeu de test s'établit à \textbf{0.009}, et l'accuracy de test atteint
\textbf{99{,}8\%}. La matrice de confusion du modèle CNN est représentée à la
figure~\ref{fig:cnn_confusion}.

\subsection{Comparaison des deux modèles}

Le tableau~\ref{tab:comparison} synthétise les principales métriques des deux modèles.

\begin{table}[h]
\centering
\label{tab:comparison}
\begin{tabular}{lcc}
\hline
\textbf{Métrique}      & \textbf{LSTM} & \textbf{CNN} \\
\hline
Test Accuracy          & 96.0\%        & 99.8\%       \\
Test Loss              & 0.127         & 0.009        \\
F1-score moyen         & 0.96          & 1.00         \\
Faux positifs          & Modéré        & Très faible  \\
Faux négatifs          & Modéré        & Très faible  \\
\hline
\end{tabular}
\caption{Comparaison des performances LSTM vs CNN}
\end{table}

\section{Analyse des résultats}

\subsection{Analyse du modèle LSTM}

Le modèle LSTM obtient une accuracy de \textbf{96\%} avec une perte de test de 0.127.
Les métriques sont équilibrées entre les deux classes, ce qui indique l'absence de biais
significatif. La matrice de confusion révèle légèrement plus de faux négatifs que de faux
positifs, comportement acceptable dans un contexte académique où les faux positifs
(accusations injustifiées) sont plus problématiques. La troncature à 512 tokens imposée
par le tokeniseur BERT constitue la principale limite de ce modèle pour les textes longs.

\subsection{Analyse du modèle CNN}

Le modèle CNN atteint une accuracy de \textbf{99{,}8\%} avec une perte de test inférieure
à 0.01, surpassant largement le LSTM. Cette performance s'explique par la capacité des
filtres convolutionnels de tailles variées (3, 4 et 5 mots) à capturer les motifs
stylistiques et lexicaux caractéristiques des textes générés par IA. La convergence dès
la première époque confirme que ces motifs sont saillants et facilement détectables.
Ces résultats sont néanmoins à nuancer : en conditions réelles, avec des textes plus
variés ou issus de modèles de génération plus récents, les performances pourraient être
inférieures.

\subsection{Comparaison de deux modéles}

Le CNN surpasse le LSTM pour trois raisons principales : la tâche repose davantage sur
la reconnaissance de motifs locaux que sur la compréhension séquentielle globale ; le CNN
converge plus rapidement ; et il est moins affecté par la troncature à 512 tokens, car il
opère sur des fenêtres locales. 

\section{Conclusion}

Ce chapitre a comparé les performances du LSTM et du CNN pour la détection de textes
générés par IA. Le CNN s'est révélé nettement supérieur (99{,}8\% contre 96\% d'accuracy,
perte quatorze fois inférieure), confirmant l'adéquation des architectures convolutionnelles
à cette tâche.